\chapter{Discussion}
\label{Chapter5}
\lhead{Chapter 5. \emph{Discussion}}

\section{Interpretation of Simulated Antenna Performance}

The simulated outputs inculcated that the dimensions were correct for resonance at 433 MHz frequency. S11 value of -21.23 dB means there is very little back propagation of power and almost 99\% of transmitted power radiated as electromagnetic waves.  This and the effective hemispherical radiation pattern confirm that this antenna is suitable for broadcasting and mass reception at base station. The addition of PLA scaffolding did cause the frequency shift which allowed the height to be compressed at a ratio of 0.92 with compared to actual height, resulting in a compactness and portability of the antenna. 

\section{Analysis of Physical Fabrication Failure}

There is a large difference between the simulation results and VNA results for the constructed antenna. The cause of resonance failure at 433 MHz is not theoretical design flow but fabrication and maintenance issue. The use of solder to affix the connection to input is non-standard for a reason. The solder joint gets shorted or if it is not fixed  

\section{Interpretation of System and Network Analysis}

Despite experiencing a fabrication failure with the Tier 2 antenna, all core objectives of the project outside of this area were successfully achieved. Of particular note is the successful implementation of the Tier 1 Civilian Node (Figure 15), which demonstrates that the core architecture consisting of an ESP32, Semtech RA02 and Bluetooth to interface with the mature Meshtastic app constitutes a robust, low overhead and immensely practical solution for the end-user component of the network. The successful range and network simulations undertaken (Figures 16 \& 17) also confirm the software and protocol aspects of the project. This corroborates that the mesh logic selected is effective, and proves that the system, when fitted with functional antennas, will be able to form a resilient, multi-hop network as intended. 

\section{Addressing Research Gaps and Limitations}

The research gaps found in previously published works were successfully represented in this work. The additionally proposed integrated two-tier system with separate civilian and official nodes is a design found nowhere else in the literature. Furthermore this work is built off the mature open-source Meshtastic developer resources and employs the legal, license-free 433MHz band. Thus, this study provides a unique and practical design overview that successfully addresses trade-offs of cost, availability and resilience. The hardware and network software for base Tier 1 node was shown to be functional, resulting in successful interaction with the smartphone application and validating the logic of the mesh protocol in simulation testing. 

The chief limitation of this work is found in its use of simulation-based verification of its most novel element. The QFH antenna proved to have excellent theoretical performance (VSWR 1.21 S11 – -21.24 dB) as evidenced by simulations but this work has not progressed to the state where realistic, extensive, real-world field testing has occurred. Therefore a physical, empirical range performance comparison has not been attempted between the standard dipole-equipped Civilian Node and the QFH-equipped Base Station results of simulation. A quantification of the difference in physical range performance results in the real-world is the most important next step in potential future work.